\SetAccent{E65B63}{FCEDEF}
\BigHero{Collaboration invitation}{LET'S BUILD YOUR INSTITUTION'S\\NEXT INNOVATION PIPELINE}{From a one-day workshop to a repeatable campus capability}
\PageOverlay{\foreach \x/\icon in {0/\faLightbulb,1/\faFlask,2/\faHeart}{\node[circle,fill=white,draw=BroGold,minimum size=11mm,text=BroGold] at ($(current page.north east)+(-17mm-14mm*\x,-25mm)$) {\icon};}}
\vspace{8mm}
\begin{BroCard}\SectionTitle{The collaboration pitch}\BodyText{Institutions rarely lack talent, ideas or ambition. The missing link is often a shared, repeatable pathway that connects academic curiosity to a problem worth solving - and then carries that problem towards an experiment, intellectual property, funding, an industry partnership or measurable student opportunity.\\[1mm]I collaborate with academic leaders to design learning engagements around the institution's real priorities. The aim is not to conduct another isolated event. It is to leave behind a common vocabulary, practical tools, visible outputs and the confidence to keep innovating.}\end{BroCard}
\CardGap
\begin{BroCard}\SectionTitle{Value-led engagement options}\JourneyRow{01}{Signature Workshop | One day}{Create a shared vocabulary, practise the frameworks and leave with participant-owned canvases, concepts and action plans.}\JourneyRow{02}{Capability Sprint | Modular series}{Combine workshops, guided application and a review clinic so faculty or student teams can apply the methods to live institutional challenges.}\JourneyRow{03}{Innovation Partnership | Custom pathway}{Integrate challenge labs, research-to-IP clinics, leadership conversations and follow-through support around a defined campus or industry priority.}\vspace{1mm}\TightText{Formats can include faculty development, student bootcamps, research clinics, industry-academia challenge labs and leadership roundtables.}\end{BroCard}
\CardGap
\begin{GoldCard}\SectionTitle{A simple path from conversation to impact}\begin{minipage}[t]{.46\linewidth}\SmallCaps{01 Listen}\par\TightText{Clarify the audience, institutional priority and desired outcome.}\par\vspace{3mm}\SmallCaps{03 Facilitate}\par\TightText{Create an energetic studio where participants build together.}\end{minipage}\hfill\begin{minipage}[t]{.46\linewidth}\SmallCaps{02 Design}\par\TightText{Tailor cases, tools, challenges and participant deliverables.}\par\vspace{3mm}\SmallCaps{04 Sustain}\par\TightText{Provide canvases and a follow-through roadmap for continued action.}\end{minipage}\end{GoldCard}
\ProgramFooter{CUSTOMISED \textbullet\ OUTCOME-LED \textbullet\ HANDS-ON \textbullet\ INSTITUTION-READY}
\newpage
\CompactHero{Institutional value}{What the institution can gain}{A collaboration is designed around visible capability, reusable tools and a practical path for continuation after the event.}
\begin{minipage}[t]{.49\textwidth}
\SideCard{What the institution can gain}{\BulletList{\item A visible portfolio of novel, relevant problems\item Stronger patent, proposal and funding concepts\item Better research stories for industry and government stakeholders\item Employer-relevant evidence of student capability\item Interdisciplinary teams aligned around shared challenges\item Reusable innovation tools that extend beyond the event}}
\CardGap
\begin{GoldCard}\SectionTitle{Value-based scoping}{\fontsize{9.2}{11.3}\selectfont\bfseries\color{BroInk}Custom programs can be discussed based on the cohort, institutional priorities and desired outcomes. Each engagement is shaped around outcomes, customisation and follow-through - not delivery hours alone.}\end{GoldCard}
\end{minipage}\hfill
\begin{minipage}[t]{.47\textwidth}
\begin{BroCard}\SectionTitle{Program portfolio}\SmallCaps{01 \textbullet\ Innovation in Academia}\par\TightText{Find consequential problems; move from incremental publication to patent, partnership and funding pathways.}\par\vspace{2mm}\SmallCaps{02 \textbullet\ The Modern-Day Professional}\par\TightText{Build judgement, narrative and inventive confidence needed to convert knowledge into opportunity.}\par\vspace{2mm}\SmallCaps{03 \textbullet\ Connecting Technology to Emotions}\par\TightText{Create human-centred technology concepts that people can trust, value and adopt.}\par\vspace{2mm}\SmallCaps{04 \textbullet\ The Art of Storytelling in Education}\par\TightText{Turn ideas, projects and research into narratives people can understand, remember and act on.}\end{BroCard}
\CardGap
\begin{DarkCard}{\color{BroGold}\fontsize{13}{15}\selectfont\bfseries Start the conversation\par}\vspace{2mm}{\color{white}\fontsize{8.8}{11}\selectfont Share your priority, participant profile and intended outcome. I will help shape an engagement relevant to your campus and designed around tangible next steps.\\[2mm]\textbf{Dr Amar Banerjee}\\+91-8237731039\\amarbanerjee23@gmail.com\\amarbanerjee23.github.io\\LinkedIn: amarbanerjee23}\end{DarkCard}
\end{minipage}
\CardGap
\begin{SoftCard}\SmallCaps{Low-risk starting path}\par\vspace{1.5mm}{\fontsize{11}{13}\selectfont\bfseries\color{BroInk}Choose one cohort. Define one measurable capability. Run the pilot. Review the evidence. Scale what works.}\end{SoftCard}
\ProgramFooter{INVITE A CONVERSATION \textbullet\ CO-DESIGN THE FORMAT \textbullet\ BUILD LASTING CAPABILITY}
\newpage
